\section{The disk occupation identity for the clock}\label{app:occupation-C16}

This appendix gives the full argument for the killed occupation identity
used in the proof of Lemma~\ref{lemma:integral-convergence-of-greens-function}.
For the circle-average clock of a $\gamma$-quantum cone, we prove the
killed occupation identity at each deterministic starting point,
including the logarithmic singularity $0$, for all $\gamma\in(0,2)$.
The coefficient is $2$ for unscaled Liouville Brownian motion and
$2/m_0$ for the process $X$ of the main text.
The latter is already the coefficient of the corrected form of
Lemma~\ref{lemma:integral-convergence-of-greens-function}.
We also give a precise killed-process version of the argument in
Lemma~\ref{lemma:integral-convergence-of-greens-function}.

\subsection{Statement, normalization, and exact inputs}

Write $B_a=B_a(0)$. Fix $\gamma\in(0,2)$, $\rho\in(0,1)$,
a deterministic disk $D=B_r(z)\subset B_\rho$, and a deterministic
$y\in D$. In particular, $\cl(D)$ is a compact subset of $B_1$.
Let $h$ be the circle-average embedding of a $\gamma$-quantum cone.
We use exactly the field and area normalization of the main text:
\begin{equation}\label{eq:occ-field}
 h|_{B_1}\overset{d}{=}(g-\gamma\log|\cdot|)|_{B_1},
 \qquad
 \mathbb{E}[g(x)g(w)]=\log\frac1{|x-w|}\quad(x,w\in B_1),
\end{equation}
where $g$ is the whole-plane GFF with $g_1(0)=0$, and
$\mu_h=\lim_{\eps\to0}\eps^{\gamma^2/2}e^{\gamma h_\eps(x)}\,dx$.
The first identity is \eqref{eq:gff} and follows directly from the
definition of the cone. In arguments
inside $B_1$ we couple the two sides of \eqref{eq:occ-field}.

Conditionally on $\mathcal F=\sigma(h,\eta)$, sample standard planar Brownian
motion $B^y$, independently of this environment. Its generator is
$\tfrac12\Delta$. Put
\begin{equation}\label{eq:occ-clock}
 \sigma=\inf\{s\geq0:B^y_s\notin D\},\qquad
 \phi_\eps(s)=\int_0^{s\wedge\sigma}
       \eps^{\gamma^2/2}e^{\gamma h_\eps(B^y_u)}\,du .
\end{equation}
Here and below the continuum Dirichlet Green function has normalization
$\Delta G_D(y,\cdot)=-\delta_y$. In a disk, if
$a=(y-z)/r$ and $b=(x-z)/r$, it is
\begin{equation}\label{eq:occ-green}
 G_D(y,x)=\frac1{2\pi}
     \log\left|\frac{1-\overline a b}{a-b}\right|.
\end{equation}
Consequently the ordinary Brownian occupation kernel is $2G_D$.
For example, this follows by applying It\^o's formula to the
zero-boundary solution of $-\tfrac12\Delta u=f$, first for smooth $f$,
and then by approximation and monotone convergence.

\begin{theorem}\label{thm:occ-main}
The clocks in \eqref{eq:occ-clock} converge, uniformly in $s\geq0$ in joint
probability over the field and Brownian motion, to a clock $\phi$.
Almost surely in the environment, $\phi$ is Brownian-almost surely
continuous, starts at $0$, is strictly increasing before $\sigma$,
and has $0<\phi(\sigma)<\infty$.

Define the unscaled process $Y_t^y=B^y_{\phi^{-1}(t)}$ before its exit
time $\tau_D^Y=\phi(\sigma)$, and the process $X$ of the main text by
$X_t^y=Y^y_{m_0t}$, with $\tau_D^X=\phi(\sigma)/m_0$.
Here $m_0$ is the unique deterministic annealed median of the raw exit time
of $Y^0$ from $B_{1/2}$, so $0<m_0<\infty$.
Almost surely in the environment, simultaneously for every bounded
Borel function $f:D\to\R$,
\begin{align}
 \mathbb{E}\left[\int_0^{\tau_D^Y}f(Y_t^y)\,dt\,\middle|\,h,\eta\right]
   &=2\int_DG_D(y,x)f(x)\,\mu_h(dx), \label{eq:occ-raw}\\
 \mathbb{E}\left[\int_0^{\tau_D^X}f(X_t^y)\,dt\,\middle|\,h,\eta\right]
   &=\frac2{m_0}\int_DG_D(y,x)f(x)\,\mu_h(dx).
       \label{eq:occ-scaled}
\end{align}
Both lifetimes have finite quenched mean. The event in this statement
may depend on the fixed parameters $\gamma,\rho,D,y$; it does not
depend on $f$. In particular the conclusion applies to every bounded
H\"older function which is measurable with respect to $(h,\eta)$.
\end{theorem}

\begin{remark}[What is, and is not, an approximation claim]\label{rem:occ-mode}
The display defining $\phi$ in the main text does not specify a mode of
convergence. Theorem~\ref{thm:occ-main} constructs its precise
circle-average approximants in the full-$\eps$ probability sense.
Every almost sure version of that displayed limit has the same value;
an almost surely convergent deterministic subsequence also exists.
We do not claim almost sure convergence along every real
$\eps\downarrow0$ at the singular start. This precision does not change
the resulting clock or its quenched occupation identity.
\end{remark}

For clarity, the exact non-elementary inputs to the proof are listed
here. Page numbers refer to the versions specified in the bibliography.
\begin{enumerate}[label=(\roman*)]
\item Berestycki \cite[Theorem 1.1, p.~3; hypotheses (1.2)--(1.5),
pp.~2--3]{berestycki2017elementary}: for covariance
$\log|x-w|^{-1}$ plus a continuous function, and a finite reference
measure of finite $(d-\delta)$-energy for every $\delta>0$, the
variance-normalized GMC approximants converge in $L^1$ and probability
when $\gamma<\sqrt{2d}$, and weakly in probability as measures.
The uniform circle mollifier is explicitly admissible on p.~3.
\item Berestycki \cite[Theorem 1.2, p.~949, and its proof, p.~959;
Remark 3.5, p.~959]{berestycki2015diffusion}: the circle-average construction for a
zero-boundary GFF converges along dyadic scales and its pre-exit clock
is continuous and strictly increasing. Strict increase is proved in
the paragraph immediately preceding Remark 3.5; that remark proves
the absence of clock jumps. We use these precise clock conclusions.
\item Berestycki--Gwynne \cite[Lemma A.1, p.~66, equation (A.1)]{berestycki2020random}:
for the circle-average $\gamma$-cone and any small $\zeta>0$, with
polynomially high probability as $\delta\downarrow0$,
\begin{equation}\label{eq:occ-BGmass}
 \delta^{(2+\gamma)^2/2+\zeta}
 \leq\mu_h(B_\delta(v))
 \leq\delta^{(2-\gamma)^2/2-\zeta},
 \qquad v\in B_{1-\delta}.
\end{equation}
\end{enumerate}
The additional exact inputs used only for Lemma~\ref{lemma:integral-convergence-of-greens-function} are recorded in
Section~\ref{sec:occ-use}. In particular we do not use the assertion of the
singular-start Green formula in \cite[footnote 2, p.~9]{berestycki2020random} as a proof
of that formula.

\subsection{The nonsingular circle-average clock}

Let
\[
 \nu_\eps(dx)=\eps^{\gamma^2/2}e^{\gamma g_\eps(x)}\,dx,
 \qquad \nu=\mu_g,\qquad
 A_\eps(s)=\int_0^{s\wedge\sigma}
                \eps^{\gamma^2/2}e^{\gamma g_\eps(B^y_u)}\,du.
\]
For sufficiently small $\eps$, all averaging circles centered in
$\cl(D)$ lie in $B_1$. Averaging the logarithmic covariance twice gives
\begin{equation}\label{eq:occ-var}
 \operatorname{Var}(g_\eps(x))=\log(1/\eps),\qquad
 \mathbb{E}\,\nu_\eps(dx)=dx.
\end{equation}
Thus there is no missing conformal-radius multiplier in this
normalization.

\begin{lemma}\label{lem:occ-base}
The clocks $A_\eps$ converge uniformly in $s\geq0$ in probability to a
continuous clock $A$, strictly increasing before $\sigma$.
For this specified deterministic $y$, on one event of full field
probability,
\begin{equation}\label{eq:occ-base}
 \mathbb{E}_y\left[\int_0^\sigma q(B^y_s)\,dA(s)\,\middle|\,g\right]
       =2\int_DG_D(y,x)q(x)\,\nu(dx)
\end{equation}
simultaneously for every bounded Borel $q$.
\end{lemma}

\begin{proof}
First verify the hypothesis needed to apply \cite[Theorem 1.1,
p.~3]{berestycki2017elementary} conditionally on a Brownian path. For $0<a<2$ and an
integer $n$, Brownian Gaussian increments give
\[
 \mathbb{E}_y\int_0^n\int_0^n |B^y_s-B^y_t|^{-a}\,ds\,dt
   =C_a\int_0^n\int_0^n |s-t|^{-a/2}\,ds\,dt<\infty.
\]
Intersecting over integers $n$ and a countable sequence of exponents
tending to $2$, and using $\sigma<\infty$ almost surely, proves that
the Brownian occupation measure
$\lambda_B(E)=\int_0^\sigma\ones_{\{B^y_s\in E\}}\,ds$
has finite $a$-energy for every $a<2$. Its support is in $\cl(D)$,
strictly inside a larger disk in $B_1$.

For this reference measure, the theorem and \eqref{eq:occ-var} give
convergence of $A_\eps(\sigma)$ in $L^1$ in the field, conditionally
on almost every Brownian path. The same applies with $\sigma$ replaced
by $t\wedge\sigma$ and with bounded nonnegative continuous spatial
weights. Each conditional $L^1$ error is bounded by twice the reference
mass, because the approximants and limit have that mass as their
field expectation. Since
$\mathbb{E}_y\sigma=(r^2-|y-z|^2)/2<\infty$, dominated convergence over Brownian
paths gives joint $L^1$ convergence. Signed bounded weights follow by
splitting positive and negative parts.

Here is an explicit transfer of clock regularity from the precise
Dirichlet-field statement in \cite[Theorem 1.2, p.~949, proof and
Remark 3.5, p.~959]{berestycki2015diffusion}. Choose $\rho<R<1$. On $B_R$, the
logarithmic covariance is the sum of the zero-boundary disk GFF
covariance
\[
 \log\frac{R|1-x\overline w/R^2|}{|x-w|}
\]
and the covariance
\[
 K_H(x,w)=-\log R-\log|1-x\overline w/R^2|
   =-\log R+
       \operatorname{Re}\sum_{n\geq1}\frac{(x\overline w/R^2)^n}{n}.
\]
The latter is positive definite: it is the covariance of an independent
real Gaussian constant of variance $-\log R$ plus a harmonic Gaussian
power series. The series converges locally uniformly almost surely.
Thus $g$ has the law of $g^{\,0}+H$ there, with independent
zero-boundary GFF $g^{\,0}$ and harmonic continuous $H$.
On $\cl(D)$, $H$ is bounded and its sufficiently small circle averages
equal $H$. The clocks are therefore related by the continuous positive
multiplier $e^{\gamma H(B^y_s)}$. The cited Dirichlet clock is continuous
and strictly increasing, and its dyadic limit agrees with the
probability limit just obtained. This constructs $A$ with those
properties. Convergence at finitely many time points, monotonicity of
$A_\eps$, and continuity of $A$ imply uniform convergence in probability
on each bounded time interval. Finally $\mathbb{P}_y(\sigma>n)\to0$ gives the
asserted uniform convergence on $s\geq0$.

For $q\in C(\cl(D))$, the ordinary Brownian Green identity gives
\begin{equation}\label{eq:occ-baseapprox}
 \mathbb{E}_y\left[\int_0^\sigma q(B^y_s)\,dA_\eps(s)\,\middle|\,g\right]
   =2\int_DG_D(y,x)q(x)\,\nu_\eps(dx).
\end{equation}
The weighted clocks on the left converge in joint $L^1$ by the preceding
conditional-GMC argument. Their probability limit is the Stieltjes
integral against $A$, by its continuity and uniform clock convergence.
Conditional expectation is an $L^1$ contraction.

The right side of \eqref{eq:occ-baseapprox} converges in $L^1$ in the field.
Indeed, first replace $G_D(y,\cdot)$ by its minimum with $M$ and apply
\cite[Theorem 1.1, p.~3]{berestycki2017elementary} with a bounded weighted Lebesgue
reference measure. The expected error, for both $\nu_\eps$ and $\nu$,
is at most
\[
 2\|q\|_\infty
       \int_{\{G_D(y,x)>M\}}G_D(y,x)\,dx\longrightarrow0.
\]
Here $\mathbb{E}\nu(dx)=dx$ follows from the same $L^1$ theorem; the logarithmic
singularity is Lebesgue integrable. This proves \eqref{eq:occ-base}
for each fixed continuous $q$.

Take a countable uniformly dense collection in $C(\cl(D))$ containing
$1$, for example the rational polynomials in the two coordinates.
On one field event, both sides define finite Borel measures on
$\cl(D)$ with identical integrals against this collection. Their
integrals agree for all continuous functions by uniform approximation,
so the measures agree. Both give zero mass to $\partial D$:
$A$ has no terminal atom and $G_D$ vanishes there.
Equality of the measures proves the assertion for every bounded Borel
$q$. This construction fixed $y$ at the outset and therefore has no
exceptional starting-point issue.
\end{proof}

\subsection{The logarithmic insertion and its actual regularization}

Choose a buffered radius $\rho<\rho_1<1$ and any
$b\in(0,(2-\gamma)^2/2)$. Apply \eqref{eq:occ-BGmass} at dyadic radii and
use Borel--Cantelli and monotonicity. Almost surely there are finite
random constants $C$ and $r_0>0$ such that
\begin{equation}\label{eq:occ-mass}
 \mu_h(B_a(v))\leq C a^b,\qquad
 v\in B_{\rho_1},\quad 0<a<r_0.
\end{equation}
On this event $\mu_h(\{0\})=0$. Off $0$, the circle average of
$-\log|\cdot|$ is $-\log|x|$ for $\eps<|x|$. Local convergence of
the area measures on compact subsets away from $0$, followed by a
countable exhaustion, therefore gives
\begin{equation}\label{eq:occ-weight}
 \mu_h(dx)=|x|^{-\gamma^2}\nu(dx)\quad\text{on }B_1.
\end{equation}
The value of the weight at $0$ is immaterial: $\nu(\{0\})=0$ follows
from $\mathbb{E}\nu(dx)=dx$, and $\mu_h(\{0\})=0$ was just proved.

Since \eqref{eq:occ-green} implies
$2G_D(v,x)\leq C_D(1+\log^+(1/|v-x|))$, a dyadic annulus sum using
\eqref{eq:occ-mass} gives
\begin{equation}\label{eq:occ-potential}
 \sup_{v\in D}\int_DG_D(v,x)\,\mu_h(dx)<\infty .
\end{equation}
For the small-distance contribution the summable upper bound is a
constant times $\sum_{j\geq0}(j+1)2^{-jb}$; the remaining contribution
is bounded by local finiteness.

Using bounded truncations of $|x|^{-\gamma^2}$ in Lemma~\ref{lem:occ-base},
define
\begin{equation}\label{eq:occ-weightedclock}
 \Phi(s)=\int_0^{s\wedge\sigma}|B^y_u|^{-\gamma^2}\,dA(u).
\end{equation}
Monotone convergence and \eqref{eq:occ-weight}--\eqref{eq:occ-potential} give
\begin{equation}\label{eq:occ-weightedmean}
 \mathbb{E}_y[\Phi(\sigma)\mid g]
       =2\int_DG_D(y,x)\,\mu_h(dx)<\infty .
\end{equation}
In particular $\Phi(\sigma)<\infty$ Brownian-almost surely.
Lemma~\ref{lem:occ-base} with $q=\ones_{\{0\}}$ shows that
$dA$ gives no mass to $\{s:B^y_s=0\}$, even when $y=0$.
One may set the weight to any positive finite value at $0$.
The finite measure defining $d\Phi$ is absolutely continuous with
respect to the atomless measure $dA$. Thus $\Phi$ is continuous,
$\Phi(0)=0$, and it is strictly increasing before $\sigma$.
In particular there is no initial holding interval or clock atom at $0$.

It remains essential to identify \eqref{eq:occ-weightedclock} with
\emph{the circle-average regularization in the main text}. Finite potentials
or a generic additive-functional correspondence alone would not do this.
We provide that identification next.

\begin{lemma}[Regularized origin tails]\label{lem:occ-tails}
Put
\[
 \mu_\eps(dx)=\eps^{\gamma^2/2}e^{\gamma h_\eps(x)}\,dx,\qquad
 0<p<\min\{1,4/\gamma^2-1\},\qquad
 \kappa=(2-\gamma^2/2)p-\gamma^2p^2/2>0 .
\]
For all sufficiently small $a,\eps>0$, uniformly in $\eps$,
\begin{equation}\label{eq:occ-tail}
 \mathbb{E}\left[\left(\int_{B_a}
          (1+\log(1/|x|))\,\mu_\eps(dx)\right)^p\right]
     \leq C a^\kappa(1+\log(1/a))^p .
\end{equation}
\end{lemma}

\begin{proof}
The circle mean-value formula, including at $x=0$, gives exactly
\begin{equation}\label{eq:occ-regweight}
 \mu_\eps(dx)=(|x|\vee\eps)^{-\gamma^2}\nu_\eps(dx).
\end{equation}
For $R<1$, the covariance in \eqref{eq:occ-field} gives on $B_1$
\[
 g(Ru)=N_R+g^R(u),\qquad
 \operatorname{Var}(N_R)=\log(1/R),
\]
where $N_R=g_R(0)$ is independent of $g^R$ and
$g^R|_{B_1}\overset{d}{=}g|_{B_1}$.
For example $\operatorname{Cov}(g(Ru),g_R(0))=\log(1/R)$ for
$|u|<1$, and subtracting this covariance gives the claim.
Whenever the averaging circles lie in $B_R$, \eqref{eq:occ-regweight}
therefore yields
\begin{equation}\label{eq:occ-scale}
 \mu_\eps(R\,du)
   =R^{2-\gamma^2/2}e^{\gamma N_R}
       (|u|\vee(\eps/R))^{-\gamma^2}\nu^R_{\eps/R}(du),
\end{equation}
where $\nu^R_\delta$ is the nonsingular regularization for $g^R$.
It has expectation $du$ on the sets used below, by \eqref{eq:occ-var}.

First let $A_s=\{s/2<|x|\leq s\}$ and $\eps\leq s/4$, and take
$R=2s$. In rescaled coordinates, $1/4<|u|\leq1/2$ and
$\eps/R\leq1/8$. The insertion weight is bounded, while the
logarithmic factor is at most $C(1+\log(1/s))$.
Independence in \eqref{eq:occ-scale}, followed by
$\mathbb{E}[V^p]\leq(\mathbb{E} V)^p$, gives
\[
 \mathbb{E}\left[\left(\int_{A_s}(1+\log(1/|x|))\,\mu_\eps(dx)\right)^p\right]
  \leq C s^\kappa(1+\log(1/s))^p,
\]
since $\mathbb{E} e^{p\gamma N_R}=R^{-\gamma^2p^2/2}$.

For the regularization core $B_c$ with $c\leq4\eps$, instead take
$R=8\eps<1$. Then $\eps/R=1/8$ and $c/R\leq1/2$.
The expected rescaled log-weighted integral is bounded by
$C(c/R)^2(1+\log(1/c))$, as follows by integrating the logarithm
against $du$; the insertion is bounded at this fixed rescaled
regularization. Thus
\[
 \mathbb{E}\left[\left(\int_{B_c}(1+\log(1/|x|))\,\mu_\eps(dx)\right)^p\right]
 \leq C R^\kappa(c/R)^{2p}(1+\log(1/c))^p
 \leq C c^\kappa(1+\log(1/c))^p,
\]
where the last inequality uses $\kappa<2p$ and $c\leq R$.
This also covers $a\leq4\eps$. Otherwise split $B_a$ into dyadic
annuli until the remaining core has radius at most $4\eps$.
Use $(\sum_j v_j)^p\leq\sum_j v_j^p$ and
\[
 \sum_{j\geq0}(2^{-j}a)^\kappa
       (1+\log(1/a)+j\log2)^p
       \leq C a^\kappa(1+\log(1/a))^p
\]
to obtain \eqref{eq:occ-tail}.
\end{proof}

\begin{proof}[Identification of the clock of the main text]
Let $0\leq\chi_a\leq1$ be continuous, equal to $0$ on $B_a$ and $1$
off $B_{2a}$.
For $\eps<a$, \eqref{eq:occ-regweight} implies
\[
 \int_0^{s\wedge\sigma}\chi_a(B^y_u)
       \eps^{\gamma^2/2}e^{\gamma h_\eps(B^y_u)}\,du
   =\int_0^{s\wedge\sigma}
        \chi_a(B^y_u)|B^y_u|^{-\gamma^2}\,dA_\eps(u).
\]
For fixed $a$ these integrals converge uniformly in $s$ in probability
to the corresponding integral against $A$. To justify this
Stieltjes continuity assertion directly, restrict first to
$\{\sigma\leq n\}$, approximate the continuous time function
$\chi_a(B^y_u)|B^y_u|^{-\gamma^2}$ uniformly by step functions,
and use uniform convergence of the integrators and their total masses.
Then let $n\to\infty$.

The removed nonnegative part is bounded, uniformly in $s$, by
\[
 R_{\eps,a}=\int_0^\sigma\ones_{\{|B^y_u|<2a\}}
                  \eps^{\gamma^2/2}e^{\gamma h_\eps(B^y_u)}\,du.
\]
Ordinary Brownian occupation gives
\[
 U_{\eps,a}(g):=\mathbb{E}_y[R_{\eps,a}\mid g]
       =2\int_{D\cap B_{2a}}G_D(y,x)\,\mu_\eps(dx).
\]
For $y=0$, $G_D(0,x)\leq C_D(1+\log(1/|x|))$ near $0$.
For $y\ne0$, the kernel is bounded there if $2a<|y|/2$.
Consequently Lemma~\ref{lem:occ-tails} and conditional Markov inequality
give, for every $u>0$,
\begin{align*}
 \mathbb{P}(R_{\eps,a}>u)
 &\leq\mathbb{E}\min\{1,U_{\eps,a}/u\}\\
 &\leq u^{-p}\mathbb{E}[U_{\eps,a}^p]
 \leq C_{D,y}u^{-p}a^\kappa(1+\log(1/a))^p\longrightarrow0 ,
\end{align*}
uniformly in sufficiently small $\eps$. The removed part for $\Phi$
tends to zero almost surely: its total mass is finite by
\eqref{eq:occ-weightedmean}, and the limiting set $\{B^y_u=0\}$ has
zero clock mass. First let $\eps\to0$ with $a$ fixed, and then
$a\to0$. We conclude that
\[
 \sup_{s\geq0}|\phi_\eps(s)-\Phi(s)|\longrightarrow0
       \quad\text{in probability}.
\]
This proves the required identification $\phi=\Phi$, including
at $y=0$, throughout $\gamma\in(0,2)$.
No annealed first moment of the inserted clock was used.
These constructions agree on nested disks up to the smaller Brownian
exit time: their stopped approximants are identical, so uniqueness
of the probability limit gives the same stopped clock. Thus this
identification also supplies the restriction of the clock of the main text
when a larger ambient disk is used.
\end{proof}

\subsection{Occupation, simultaneity, and the coefficient}

\begin{proof}[Proof of Theorem~\ref{thm:occ-main}]
The clock construction and its regularity have been proved above.
For a fixed environment, extend \eqref{eq:occ-base} by monotone convergence
to nonnegative Borel functions. Equations \eqref{eq:occ-weight} and
\eqref{eq:occ-weightedclock} then identify, as finite Borel measures,
\[
 \mathbb{E}_y\left[\int_0^\sigma\ones_{\{B^y_s\in dx\}}\,d\phi(s)
                 \,\middle|\,g\right]
     =2G_D(y,x)\,\mu_h(dx).
\]
The equality holds on the same event as in Lemma~\ref{lem:occ-base}
and \eqref{eq:occ-potential}, simultaneously for all Borel sets.
In particular its integral against any bounded Borel $f$ is legitimate
and absolutely integrable.

A continuous strictly increasing clock maps $[0,\sigma]$ onto
$[0,\phi(\sigma)]$. The change of variables for its Stieltjes measure
therefore gives, pathwise,
\[
 \int_0^{\tau_D^Y}f(Y_t^y)\,dt
       =\int_0^\sigma f(B^y_s)\,d\phi(s).
\]
This proves \eqref{eq:occ-raw}. Strict positivity and finiteness of the
raw exit time also show that its annealed median from $0$ in $B_{1/2}$
is deterministic, positive, and finite. To check uniqueness, take
$1/2<R<1$ in the Gaussian decomposition used in
\eqref{eq:occ-scale}. The raw exit time $L$ has the representation
$L=e^{\gamma N_R}L'$, with $N_R$ a nondegenerate Gaussian independent
of $0<L'<\infty$. Hence $\log L$ has a strictly positive continuous
density, obtained by convolving the law of $\log L'$ with a Gaussian
density. Its median is unique. Finally $X_t=Y_{m_0t}$,
so changing variables $s=m_0t$ proves \eqref{eq:occ-scaled}.

The argument is local in $h|_{B_1}$, so the law identity
\eqref{eq:occ-field} transfers it to the cone. Conditional Brownian
sampling is independent of the entire pair $(h,\eta)$; all kernels
used above are functions of $h$ alone. Conditioning additionally on
$\eta$ therefore leaves them unchanged. The finite-measure identity
holds before choosing $f$, which proves the claimed simultaneity and
allows environment-measurable $f$. Complex-valued tests follow by
treating their real and imaginary parts separately.
\end{proof}

The coefficient of the named process $X$ is therefore $2/m_0=\lambda_0$, the constant of the main text;
coefficient $2$ is the coefficient of the raw process $Y$.

\begin{remark}\label{rem:occ-m0}
Here $m_0$ must be defined with the starting point $0$:
\[
 m_0=\operatorname{median}
       \left(\inf\{t\geq0:|Y_t^0|\geq1/2\}\right).
\]
It cannot define a different $m_0$ for each displayed $z$:
the discrete median in \eqref{eq:meps} is taken from $0$,
and a start outside $B_{1/2}$ would give median zero.
\end{remark}

\subsection{The use in Lemma~\ref{lemma:integral-convergence-of-greens-function}}\label{sec:occ-use}

The conclusion of Lemma~\ref{lemma:integral-convergence-of-greens-function}
needs no change of coefficient or quantifiers. Its proof should explicitly
use killed occupation in $D=B_r(z)$. A path frozen at its boundary point
has infinite occupation of $f=1$ over $[0,\infty)$; stopping only on
leaving $B_\rho$ does not give the Green function of $D$.

The following argument gives the convergence-to-occupation step of the
proof of Lemma~\ref{lemma:integral-convergence-of-greens-function}.
For the embedded walk, put
\[
 I_\eps=\VGeps(D)\setminus\VGeps(\partial D),\qquad
 K_\eps=\inf\{k\geq0:X_k^{y,\eps}\notin I_\eps\},\qquad
 W^\eps(t)=\eta(X^{y,\eps}_{\meps t}),
\]
with linear interpolation in the last expression. Use full walks,
or walks observed until leaving a larger buffered disk.
Let $\tau=\tau_D^X$ for the accelerated continuum process.

The exact external inputs here are
\cite[Theorem 1.2, p.~10]{berestycki2020random}, giving quenched convergence in
probability in the topology of local uniform convergence of paths;
\cite[Remark 1.3, equation (1.8), p.~10]{berestycki2020random}, giving
$m_\eps\asymp\eps^{-1}$; \cite[Lemma 2.4, equation (2.8),
p.~17]{berestycki2020random}, giving vanishing maximal cell diameter on buffered
compact disks with polynomially high probability; and
\cite[Proposition 5.1, equation (5.1), p.~28]{berestycki2020random}, giving,
on an event of polynomially high probability, the bounds
\begin{equation}\label{eq:occ-discretemom}
 \mathbb{E}[(\tau^\eps_U)^N\mid h,\eta]
 \leq N!C^N\eps^{-N}
  \left(\sup_{v\in U}
       \int_{B_{\eps^\alpha}(U)}
         \left(\log\frac1{|v-w|}+1\right)\mu_h(dw)\right)^N
\end{equation}
for all positive integers $N$ and admissible Borel
$U\subset B_{\rho_1}$ with $\mu_h(U)\geq\eps^\alpha$.
Its $\tau_U^\eps$ is exit from the cells intersecting $U$;
$K_\eps$ kills no later, so it satisfies the same upper bound.

Start with any deterministic sequence $\eps\downarrow0$.
Pass to a deterministic subsequence on
which the quenched path laws converge almost surely, the maximal
relevant cell diameter $d_\eps$ tends to zero almost surely, and
\eqref{eq:occ-discretemom} holds eventually. The mass condition for the
fixed disk holds eventually because $\mu_h(D)>0$ almost surely.
The logarithmic integral in \eqref{eq:occ-discretemom} is bounded uniformly
on its tail by the same annulus argument as \eqref{eq:occ-potential},
using the buffered radius and \eqref{eq:occ-mass}. Taking $N=2$ yields
a finite random $M$ with
\begin{equation}\label{eq:occ-UI}
 \mathbb{E}[(K_\eps/m_\eps)^2\mid h,\eta]\leq M,\qquad
 \mathbb{E}[(K_\eps/m_\eps-T)_+\mid h,\eta]\leq M/T .
\end{equation}

Fix an environment in this event and in the event of
Theorem~\ref{thm:occ-main}, also applied in a buffered disk. Couple its
quenched path laws by the Skorokhod representation, so that
$W^\eps\to X$ locally uniformly almost surely. Then
\begin{equation}\label{eq:occ-exitconv}
 K_\eps/m_\eps\longrightarrow\tau
       \quad\text{almost surely under the coupling}.
\end{equation}
We give the pathwise detail since exit time is not continuous at
every continuous path. If $t<\tau$, the compact path $X([0,t])$
has positive distance from $D^c$. Uniform convergence and shrinking
cells imply that every visited cell before time $m_\eps t$ lies
inside $D$ and misses $\partial D$. This proves the lower bound
in \eqref{eq:occ-exitconv}.

At ordinary Brownian exit from a disk, the projection of subsequent
increments onto the outward normal is one-dimensional Brownian
motion by the strong Markov property. The reflection principle
shows it has positive values arbitrarily soon. Thus Brownian motion
visits the strict exterior of the closed disk arbitrarily soon after
exit. The continuous strictly increasing clock in the buffered disk
transfers these visits to $X$ at times arbitrarily close after
$\tau$. At any such time, the exterior point is at positive distance
from $\cl(D)$, and this persists under uniform convergence.
A linear segment with a point outside the convex closed disk has
an endpoint outside it; that endpoint's cell cannot belong to
$I_\eps$. The corresponding mesh time differs by at most
$m_\eps^{-1}\to0$. This proves the upper bound and hence
\eqref{eq:occ-exitconv}.

Let $f$ be any bounded H\"older function on $D$, measurable with
respect to the environment. It has a unique continuous extension
to $\cl(D)$. Compose that extension with nearest-point projection
onto the closed disk to obtain a bounded H\"older extension
$\widetilde f$ to $\C$. After the environment is fixed this is
an ordinary fixed function, of some exponent $\theta>0$ and
constant $H<\infty$. The coupling gives, for each fixed $T$,
\[
 \int_0^{T\wedge(K_\eps/m_\eps)}
             \widetilde f(W^\eps(t))\,dt
       \longrightarrow \int_0^{T\wedge\tau}f(X_t)\,dt .
\]
The bound $T\|f\|_\infty$ permits passage to quenched expectations.
Adjacent cells intersect, so the relevant interpolation segments
have length at most $2d_\eps$. The difference from the discrete
Riemann sum over $0\leq k<K_\eps$, $k<m_\eps T$, divided by
$m_\eps$, is at most
\[
 T H(2d_\eps)^\theta+2\|f\|_\infty/m_\eps .
\]
The time horizon is fixed here; a uniform error independent of all
$T>0$ is unnecessary.

By \eqref{eq:occ-UI}, the expected discarded discrete tail is at most
$\|f\|_\infty(M/T+m_\eps^{-1})$. The continuum tail tends to zero
because $\mathbb{E}[\tau\mid h,\eta]<\infty$ by
Theorem~\ref{thm:occ-main}. Letting first $\eps\to0$ and then $T\to\infty$
proves
\begin{align*}
 m_\eps^{-1}\sum_{b\in I_\eps}
       \Gr{{\eps}}{D}(a_y^{{\eps}},b)f(\eta(b))
 &=m_\eps^{-1}\mathbb{E}\left[\sum_{k<K_\eps}
             f(\eta(X_k^{y,\eps}))\,\middle|\,h,\eta\right]\\
 &\longrightarrow
   \mathbb{E}\left[\int_0^\tau f(X_t^y)\,dt\,\middle|\,h,\eta\right]\\
 &=\frac2{m_0}\int_DG_D(y,x)f(x)\,\mu_h(dx).
\end{align*}
The sum may equivalently range over all cells intersecting $D$
if one uses $\widetilde f$: their killed Green function is zero
outside $I_\eps$. This avoids evaluating $f$ at a representative
outside its domain.

Every deterministic input sequence has a deterministic subsequence
with this almost sure conclusion. The subsequence criterion gives
convergence in probability along all meshes, as required in
Lemma~\ref{lemma:integral-convergence-of-greens-function}. The fixed-environment argument works for all bounded
H\"older tests on the same event, so it also justifies the
environment-dependent tests in that lemma.
